\subsection{Second-type cancellation by reversible lattice moves}
\label{sec:flag-second-type-cancellation}

Keep the actual local data, notation and rank range of
Sections~\ref{sec:flag-generic-local-amplitude} and
\ref{sec:flag-defect-stratification}. Write \(R=O_{F_v}\),
\(h(v,w)=\langle v,w\rangle\) (linear in the first variable),
\(D\) for the joint frame span (called \(J\) in (GL1)),
\(W_0=D^\perp\), and
\[
 \rho=n-4s=k-2r-3\ge4,\qquad \epsilon=q^{1-\rho}.
\]
The two adapted frames are denoted \(f_i^*=u_i\) and
\(f_i^{\prime *}=u_i'\), \(0\le i\le r\).
The character \(\psi_0\) is the original additive character
of \(E_v\), with annihilator of \(O\) under multiplication
equal to \(O\); on the base field we use a primitive trace
character of conductor \(R\). No odd residue characteristic
assumption is made.

\begin{lemma}[Uniform quadratic Fourier and leakage estimates]
\label{lem:flag-second-type-fourier}
Let \(\Lambda=\Lambda^\#\), \(Y\subset D\) be any subspace,
\(U=\Lambda\cap Y^\perp\), \(W=\Lambda\cap W_0\).
Use probability Haar measure on \(U\), put \(Q(v)=h(v,v)\),
and let \(Z_a=\{v\in U:Q(v)\in\pi^a R\}\), \(a\ge0\).
Then
\[
 \left|q^a\Pr(Z_a)-1\right|\le\epsilon,\qquad
 q^a\left|\int_{Z_a}\chi(v)\,dv\right|\le\epsilon
 \quad(\chi\text{ a nontrivial continuous unitary character of }U).
 \tag{SC1}
\]
Moreover
\[
 q^a\Pr\{v\in Z_a:\nexists w\in W,\ h(v,w)=1\}
 \le 2(1+\epsilon)\epsilon .                         \tag{SC2}
\]
In the nonsplit case the right side of (SC2) can be replaced
by \((1+\epsilon)\epsilon^2\).
\end{lemma}
\begin{proof}
The saturated lattice \(W\) has rank \(n-2s\).
Its reduction embeds in the perfect hermitian space
\(\Lambda/\pi\Lambda\), so the radical of its restriction
has dimension at most \(2s\). In the split case use the
two component subspaces of dimension \(n-2s\) and their
pairing; its rank is again at least \(n-4s\).
Choose a nondegenerate subspace of rank \(\rho\) (in the
split case paired subspaces of that rank), and lift its
basis to \(W\). The resulting unimodular hermitian
sublattice has an orthonormal basis: over the unramified
field split off a unit-norm line successively, rescaling
by a unit of the prescribed norm. Such units exist by
surjectivity of the residue norm and successive trace
lifting. Over the split algebra choose dual bases.
This also works at \(q=2\). Orthogonal projection onto
this summand is integral on \(\Lambda\). Consequently
\(U=O^\rho\perp U_1\), with the displayed summand in \(W\).

Fourier inversion on \(R/\pi^a R\) gives
\[
 q^a 1_{Z_a}(v)=\sum_{t\in\pi^{-a}R/R}\psi_F(t Q(v)).
\]
There are \(q^i-q^{i-1}\) characters of conductor exponent
\(i>0\). On one orthonormal coordinate, for any linear
unitary additive character \(\eta\) of \(O\), write
\(I=\int_O\psi_F(t z\bar z)\eta(z)\,dz\).
Squaring its modulus and substituting \(z=z'+b\) shows
that the inner integral in \(z'\) vanishes unless
\(b\in\pi^i O\), by the perfect trace pairing of the
unramified quadratic algebra (also in the split case).
Thus \(|I|^2\le q^{-2i}\), uniformly in \(\eta\).
Factoring the integral on the orthogonal decomposition,
and bounding the \(U_1\)-factor by one, gives \(q^{-i\rho}\).
The zero-frequency term is one for the trivial character
and zero for a nontrivial character. The total nonzero
frequency error is at most
\[
 \sum_{i\ge1}(q^i-q^{i-1})q^{-i\rho}
 =\frac{(1-q^{-1})\epsilon}{1-\epsilon}\le\epsilon .
\]
This proves (SC1), including \(a=0\).

Failure of the pairing condition implies that the image
ideal \(h(v,W)\subset O\) is proper. In particular all
\(\rho\) orthonormal coordinates lie in one maximal ideal
\(\mathfrak p\) of \(O\). Write it as \(\alpha O\),
where \(\alpha=\pi\) in the nonsplit case and
\(\alpha=(\pi,1)\) or \((1,\pi)\) in the split case.
Then \([O:\alpha O]=q^b\), \(N(\alpha)=\pi^b\),
with \(b=2\) or \(1\), respectively. The restriction
of the \(\rho\) coordinates costs \(q^{-b\rho}\);
after rescaling their quadratic form is multiplied
by \(\pi^b\). Fourier inversion, bounding the
frequencies \(i\le b\) trivially, gives
\[
 q^a\Pr(Z_a,\ z_1,\dots,z_\rho\in\alpha O)
 \le q^{-b\rho}\left(q^b+
     \sum_{i>b}(q^i-q^{i-1})q^{-(i-b)\rho}\right)
 \le (1+\epsilon)q^{b(1-\rho)} .
\]
The union of the two maximal ideals in the split case,
or the single ideal in the nonsplit case, proves (SC2).
\end{proof}

\begin{lemma}[The actual reversible moves]
\label{lem:flag-second-type-moves}
Fix a datum \(\delta\) and a nonempty original
\(\mathrm{InFF}(\beta,j,\beta_j)\) family, where
\(j\) is the largest defect level and \(\beta_j\) the
number of classes of that level. Put
\(X_j=\operatorname{span}(f_i:i<j)\),
\(X'_j=\operatorname{span}(f'_i:i<j)\), \(Y=X_j+X'_j\).
For a term \(\Lambda\) in the family let
\[
 B_\Lambda=\{(x,x')\in X_j\times X'_j:
                          x+x'\in\theta\Lambda\},\qquad
 U_\Lambda=\Lambda\cap Y^\perp .
\]
Take a pair \(b=(x,x')\in B_\Lambda\) of exact level
\(j\) (the level of \(\theta^{-1}(x+x')\)),
\(v\in U_\Lambda\), and \(\sigma\in\pi^aO\) with
\(\sigma+\bar\sigma=h(v,v)\).
Define
\[
 e(z)=z+h(z,v)x-h(z,x)(v+\sigma x),\qquad
 \Lambda_1=e^{-1}\Lambda .
 \tag{SC3}
\]
Then \(\Lambda_1\) is a term over the same datum;
the phase changes by
\[
 \omega(\Lambda_1)=\omega(\Lambda)\chi_b(v),\qquad
 \chi_b(v)=\psi_0\left(-\theta^{-1}
                h(cu_j+c'u'_j,v)\right),           \tag{SC4}
\]
where \(c,c'\) are the coefficients at index \(j-1\)
in \(x,x'\) in their respective adapted frames.
For every exact-level pair \(b\), \(\chi_b\) is
nontrivial on \(U_\Lambda\).
If \(h(v,w)=1\) for some \(w\in\Lambda\cap W_0\),
then \(\Lambda_1\) belongs to the same family.
Whenever it belongs to that family, the reverse
move is \((x,x',-v,\bar\sigma)\).
\end{lemma}
\begin{proof}
All three family parameters are witness independent
as in the term and frame-class construction.
For clarity the defining second-type conditions are:
\(\beta>1\), all defect levels at most \(r\), and
\(h(p_1,p'_2)\in\theta O\) whenever
\(p_\nu\in X,\ p'_\nu\in X'\) and
\(p_\nu+p'_\nu\in\theta\Lambda\).
Here \(X=X_r,\ X'=X'_r\). In particular \(1\le j\le r\).
In the independent adapted frame coordinates,
the level of \(y\in M_h=\Lambda\cap D\) is the
least \(d\) such that for each \(i\ge d\) the
coefficients \((A_i,B_i)\) of \(\theta y\) satisfy
\[
 B_i\in O,\quad A_i+B_i\in\theta O.
 \tag{SC5}
\]
This is exactly the source's RR condition and
is equivalent to being represented modulo \(M\)
by a vector in \(X_d+X'_d\) for \(d\le r\).
In particular level zero is equivalent to membership
in \(M\).

Put \(z_0=(x+x')/\theta\in M_h\).
By (GL3), \(\pi^a z_0\in M\); independent
coordinates in the two frames imply
\(\pi^a x,\pi^a x'\in\Lambda\).
Thus \(\sigma x,\bar\sigma x,\sigma x',\bar\sigma x'\)
are integral. We also have
\(h(x',x)=\theta C\), \(C\in O\), by the second-type
condition and conjugation. Since \(X,X'\) are
isotropic and \(v\perp Y\), expansion shows that
\(e\) is an isometry with inverse obtained by
\((v,\sigma)\mapsto(-v,\bar\sigma)\).
It fixes \(X_j\), acts by \(z\mapsto z+h(z,v)x\)
on \(X^\perp\), and preserves the flag with identity
on its successive quotients, hence belongs to \(N\).
The analogous \(e'\) built from \(x'\) belongs to \(N'\).

For \(z\in\Lambda\cap (X')^\perp\), direct substitution
gives the stronger integral congruence
\[
 ee'z-z=\theta\{h(z,v)z_0-
          (h(z,z_0)+C h(z,v))(v+\sigma x)\}
          \in\theta\Lambda .                       \tag{SC6}
\]
The same formula with inverse parameters holds on
\(\Lambda_1\cap(X')^\perp\): indeed \(x+x'\in
\theta\Lambda_1\) follows by applying \(e\) to
\(x+x'\), and \(-v\in\Lambda_1\) follows from
\(e(v)=v+(\sigma+\bar\sigma)x\).
Both \(\sigma x,\bar\sigma x\) belong to \(\Lambda_1\)
because \(e\) fixes \(x\). All integrality hypotheses
of (SC6) therefore hold on the new lattice.
These two congruences imply that \(e'\Lambda\)
and \(e^{-1}\Lambda\) have the same intersection
with \((X')^\perp\).

We record the integral extension fact used here.
If two self-dual lattices share their trace on the
orthogonal of a totally isotropic subspace \(I\),
there is an isometry fixing \(I^\perp\) pointwise
taking one to the other. To see this, choose a basis
of the common primitive intersection with \(I\).
For either lattice the pairing with that basis is
onto its integral dual, since the intersection is
a direct summand and the lattice is self-dual.
Its kernel is the common trace, and the trace modulo
the intersection has perfect induced hermitian form.
Choose the same lifts of a basis of this quotient
in the trace, giving a unimodular orthogonal summand.
Choose dual vectors to the basis in \(I\) in the
orthogonal complement of that summand in each
lattice, then modify the dual vectors by the original
isotropic basis to make their Gram matrix zero.
The required hermitian matrix is of the form
\(A+A^*\) over \(O\), by trace surjectivity on the
diagonal. The resulting hyperbolic bases and common
unimodular summand define the desired isometry.
The trace spans \(I^\perp\). This proof applies
componentwise over the split algebra and at dyadic
places.

Apply this fact with \(I=X'\), to obtain \(g\)
fixing \((X')^\perp\) pointwise and taking
\(e'\Lambda\) to \(\Lambda_1\). The new witnesses
are \(e^{-1}n\) and \(ge'n'\), where \(n,n'\)
are the old witnesses. The fixed-point property
of \(g\) implies \(g\in N'\) and \(\psi'(g)=1\).
The new frame link modulo \(\theta\Lambda_1\)
follows from (SC6) on each \(u_i'\), together
with the old link. For the graph datum take old
representatives \(z'\in\Lambda\cap(X')^\perp\),
\(z\in\Lambda\cap X^\perp\); replace them by
\(e'z'\) and \(e^{-1}z\). Their quotient classes
do not change, and by (SC6)
\[
 e'z'-e^{-1}z\in\theta\Lambda_1
      \quad\Longleftrightarrow\quad z'-z\in\theta\Lambda .
\]
The inverse congruence proves the reverse implication
for every graph pair of the new lattice. Thus the
same inner data and exact level graph are retained.

The adjacent-coefficient character of \(e\) is
\(\psi_0(\theta^{-1}c\,h(u_j,v))\).
Indeed on \(X^\perp\), \(e(f_i)=f_i+h(f_i,v)x\);
in the sum of adjacent coefficients, entries before
\(j\) vanish by \(v\perp X_j\), entries after \(j\)
vanish because \(x\in X_j\). The surviving coefficient
of \(x\) is \(c\) (unchanged by unipotent frame change),
and \(u_j-f_j\in X_j\). The primed calculation is
identical. Multiplicativity of the two characters,
and \(\psi'(g)=1\), prove (SC4).

To prove nontriviality put \(Z=cu_j+c'u'_j\).
If \(\chi_b\) is trivial on \(U_\Lambda\), scaling
by all \(O\)-multiples and the conductor property
imply \(h(\theta^{-1}Z,U_\Lambda)\subset O\).
The dual of \(U_\Lambda\) in \(V\) is \(\Lambda+Y\):
extend the induced integral functional from the
saturated direct summand \(U_\Lambda\) to \(\Lambda\),
then use self-duality; the annihilator of its span
is \(Y\). Write \(\theta^{-1}Z=z+y\), \(z\in\Lambda\),
\(y\in Y\). Then \(z\in M_h\), but the coefficient
pair at index \(j\) of \(\theta z\) is \((c,c')\).
Exact level \(j\) of \(b\) means that the pair
\((c,c')\) fails (SC5). Thus \(z\) has level at
least \(j+1\), impossible.

It remains to prove preservation of the family,
not merely of the datum. Write \(M_1\) and \(M_{h,1}\)
for the new base and saturated lattices.
For \(i<j\) both adapted vectors remain unchanged,
and the other vectors remain jointly independent.
Thus \(M_1\cap Y=M\cap Y\) and levels on \(Y\)
coincide. For \(y\in\Lambda\cap Y\), the defining
second-type test on \(x+x'\) and \(\theta y\)
implies \(h(y,x)\in O\). Since \(h(y,v)=0\),
\(e(y)=y-h(y,x)(v+\sigma x)\in\Lambda\).
Hence \(B=\Lambda\cap Y\subset\Lambda_1\cap Y\).

Suppose there is \(w\in\Lambda\cap W_0\) with
\(h(v,w)=1\). For \(y\in M_{h,1}\), one has
\(e(y)\in\Lambda\); pairing with \(w\) yields
\(h(y,x)\in O\). Therefore
\(y+h(y,v)x\in M_h\). By the level bound, write
this as \(y_0+m\) with \(y_0\in B,\ m\in M\).
The new adapted vectors are \(u_i-h(u_i,v)x\)
and \(u_i'+h(u_i',v)x'\). Write
\(m=\sum A_i u_i+\sum B_i\theta^{-1}(u_i'-u_i)\)
with integral coefficients. The element \(m_1\)
in \(M_1\) given by the same coefficients satisfies
\[
 m_1=m-h(m,v)x+
        \left(\sum B_i h(u_i',v)\right)z_0 .
\]
The last term belongs to \(B\). Also
\(h(y,v)=h(m,v)\), so \(y\in B+M_1\).
Consequently \(B\) surjects onto both \(M_h/M\)
and \(M_{h,1}/M_1\), with the same kernel and
the same levels. This proves preservation of
\(\beta,j,\beta_j\) and nonzero defect.

For the remaining second-type test, take
\(p+p'\in\theta\Lambda_1\), \(p\in X,\ p'\in X'\).
Write it as \(\theta y_0+\theta m_1\) using
\(y_0\in B\). Since the two last coefficients
vanish, the two components of \(\theta m_1\)
in \(X,X'\) are integral adapted-frame combinations
in \(\Lambda_1\). Pairing either such component
with the opposite component of any other linked pair
is in \(\theta O\), by isotropy and integrality
of the sum of that other pair divided by \(\theta\).
Thus the cross test reduces to the representatives
\(\theta y_0\in\theta\Lambda\), where it is already true.
This proves second-type preservation. Finally, if
the target is in the same family, the inverse
parameters satisfy all the original move conditions:
the link and integrality were checked above,
\(v\perp Y\), the level on \(Y\) is unchanged,
and \(\bar\sigma+\sigma=Q(-v)\).
\end{proof}

\begin{theorem}[Full second-type and nongeneric local estimates]
\label{thm:flag-second-type-cancellation}
For any fixed datum and any second-type family \(\mathcal F\),
\[
 |A_\delta(\mathcal F)|
 \le 7 q^{1-\rho}\,m_\delta(\mathcal F).              \tag{SC7}
\]
The same inequality holds for the union of all second-type
families (not for arbitrary subsets of a family).
With the original two stabilizer weights,
\[
 \|\,\mathrm{nongeneric}\,\|_{\rm amp}^2
 \le P^{39s^2}q^{-(k-2r-5)}\,nn.                     \tag{SC8}
\]
\end{theorem}
\begin{proof}
For each term in \(\mathcal F\), use the original
move coordinates \((b,v,t)\) with
\(\sigma=Q(v)/2+t\). Normalize Haar measure
on the product of the lattices \(B_\Lambda,U_\Lambda\)
and \(T=\pi^a O\cap\ker\operatorname{Tr}_{E/F}\).
The fraction of exact-level pairs is
\(\beta_j/\beta\): addition divided by \(\theta\)
identifies \(B_\Lambda\) with \(B\), which surjects
onto the defect quotient. For each \(v\) the allowable
\(t\)'s have measure one if \(v\in Z_a\), and zero
otherwise, since trace \(O\to R\) is onto and each
nonempty fibre is a translate of \(T\).
The division by 2 is only in the field, not in \(O\).

Define a nonnegative kernel \(K(\Lambda,\Lambda_1)\)
on the actual family by integrating the indicator of
the move (SC3) with target \(\Lambda_1\), multiplied
by \(q^a\beta/\beta_j\).
This is the source's \(\mathrm{kern}\): on the compact
support the finite index integral is the normalized
Haar integral of the locally constant integrand.
Indeed the pair and vector lie in compact lattices,
and \(t\) in a bounded closed subset of the trace-zero
space; the level condition is a finite coset test.
The isometry varies continuously, while the lattice
stabilizer is open. Thus the map to lattices is
locally constant and has finite image on this
compact domain. The scalar, indicator and phase
integrals can equally be evaluated on a finite
common congruence quotient. There is no convergence
assumption on an infinite family of moves.

The product Haar normalizations are identical at
every term of \(\mathcal F\). To verify this,
\(M\cap Y\) has basis the first \(j\) adapted
vectors and their linked differences divided by
\(\theta\). Its covolume is constant, because the
two flag basis changes are unipotent triangular.
Moreover \([B:M\cap Y]=\beta\); hence \(B\) and
\(B_\Lambda\) have fixed covolume. The restriction
of the hermitian pairing maps \(\Lambda\) onto
the integral dual of \(B\) in \(Y^*\), with kernel
\(U_\Lambda\), because \(B\) is a saturated direct
summand of \(\Lambda\). The quotient ambient map
\(V\to Y^*\) is fixed even when \(Y\) is degenerate.
Since the covolume of \(\Lambda=\Lambda^\#\) in \(V\)
is fixed, the covolume of \(U_\Lambda\) is fixed too.
The scalar reference \(T\) does not change.
The inverse move is \((b,v,t)\mapsto(b,-v,-t)\),
so the kernel is symmetric. It is equivariant
under the common stabilizer \(S_t=N\cap N'\):
the isometry maps both flags and all move conditions
to themselves, transports the reference lattices,
and conjugates the Eichler formula. The datum,
family and phase are invariant under \(S_t\)
(the original common-stabilizer character agreement
and witness independence used in the class unfolding).

Before discarding moves leaving \(\mathcal F\),
the row mass differs from one by at most \(\epsilon\)
by (SC1). The phase-weighted row sum has modulus
at most \(\epsilon\) by (SC4) and (SC1), since every
exact-level pair has a nontrivial character.
By the preservation lemma all discarded moves
are in the leakage set (SC2). Consequently
\[
 \left|\sum_{\Lambda_1}K(\Lambda,\Lambda_1)-1\right|,
 \quad\left|\sum_{\Lambda_1}K(\Lambda,\Lambda_1)
                           \omega(\Lambda_1)\right|
       \le \epsilon+2(1+\epsilon)\epsilon .        \tag{SC9}
\]
These sums run over actual lattices, not orbit
representatives.

For the quotient step put \(G_\Lambda=
S_t\cap\operatorname{Aut}(\Lambda)\).
The reciprocal weight \(\mu(\Lambda)\) has the
fixed normalization used in (GL6). For a pair
orbit the mass of targets at fixed source is
\[
 \mu(\Lambda)[G_\Lambda:G_\Lambda\cap G_{\Lambda_1}]
 =1/\operatorname{vol}(G_\Lambda\cap G_{\Lambda_1}),
\]
with the same reference normalization on either
side; equivalently use multiplicativity of finite
relative indices. It is unchanged on reversing
the pair. The weights and phases are orbit invariant,
and there are finitely many term orbits for fixed
datum by the local class and gluing counts (GL4)--(GL6)
and (DS5). Summing over pair orbits therefore proves
\[
 \sum_{[\Lambda]}\mu(\Lambda)\sum_{\Lambda_1}
 K(\Lambda,\Lambda_1)\omega(\Lambda_1)
 =
 \sum_{[\Lambda]}\mu(\Lambda)\omega(\Lambda)
                        \sum_{\Lambda_1}K(\Lambda,\Lambda_1).
\]
Apply the two row estimates (SC9) to this identity,
using \(|\omega|=1\). Since \(\epsilon\le1/8\), the
resulting factor \(2(3+2\epsilon)\epsilon\) is
at most \(7\epsilon\). This proves (SC7).
The nonempty families partition the second-type
terms uniquely, so the ordinary triangle inequality
at each datum gives (SC7) for their union without
any factor counting families.

The proof of (DS3) applies also to the square of
the positive mass, as it used only positive mass
bounds. Taking \(I=1\) there, the squared norm of
the second-type amplitude is bounded by
\[
 49 P^{2s^2+10s+11}q^{-1}\epsilon^2 nn .
\]
By (DS1) every other nongeneric term has strict
image, hence by (DS3) contributes squared norm at
most \(P^{2s^2+10s+11}q^{-l}nn\).
Use the squared triangle inequality with factor 2.
In any nonempty nongeneric branch \(P\ge2\);
\(49\le P^6\), \(l=k-1\ge k-2r-5\), and
\(q^{-1}\epsilon^2=q^{1-2\rho}\le
q^{-(k-2r-5)}\). The total is at most
\(P^{2s^2+10s+19}q^{-(k-2r-5)}nn\),
which implies (SC8) since
\(2s^2+10s+19\le39s^2\) for \(s\ge1\).
For \(P=1\) there are no nongeneric terms; for
\(r=0\) the second-type family is empty.
\end{proof}

The reconstruction follows the actual second-type family
of the pinned source, not an invented graph on isometry
classes. In particular (SC7) uses the positive reversible
kernel before passing to the quotient, and (SC8) uses the
original complement-datum weights. Its new ingredients
relative to (DS4) are the no-leak move preservation and
the uniform nontrivial Fourier estimate. The orthogonal
summand argument removes a conductor-dependent finite
Gauss-sum derivation; the final sum avoids the count of
second-type families. This is a finite-place supplier,
not the generic nonidentity return or the global
positive-trace estimate.
